\documentclass[11pt]{article}
\usepackage{problem_set_style}
\usepackage{kotex}
\usepackage{booktabs}
\usepackage{array}
\begin{document}
\heading{G1}{중간 · 기말 공통}{}{Spring 2027}{논문(보고서) 작성 가이드라인}

\noindent{\small 이 수업의 프로젝트 보고서는 \textbf{작은 논문}이다. 학회에 내는 논문과 같은 규칙으로 쓰고, 같은 방식(KSA Review)으로 서로 심사한다. 이 문서는 ICLR 2027 · NeurIPS 2026 · ICML 2026의 AI 정책과 The Web Conference 2027 심사 지침을 참고해 우리 수업에 맞게 만든 것이다. 학회별 비교는 리뷰 가이드라인 끝의 표를 보라.}

\section*{먼저 이것만은}
\begin{enumerate}
\item \textbf{주장은 하나, 첫 문단에.} 읽는 사람이 첫 문단만 보고 ``이 팀은 무엇을 보였다고 주장하는가''를 한 문장으로 말할 수 있어야 한다.
\item \textbf{모든 숫자는 실험에서 나와야 한다.} 표·그래프의 숫자는 제출한 코드를 실행하면 다시 나와야 한다. 지어낸 숫자, 고른 숫자, 출처를 모르는 숫자는 없어야 한다.
\item \textbf{AI는 자유롭게 써도 된다.} ``AI 사용'' 한 줄만 적는다. 대신 AI가 만든 것이 틀렸다면 그 책임은 팀에 있다 (3절).
\item \textbf{참고문헌은 직접 읽은 것만.} 전문을 열어 읽은 문헌만 인용하고, 링크와 원문 발췌를 \textbf{인용 근거 표}로 붙인다 (4절). \textbf{존재하지 않는 참고문헌은 0점}이다.
\item \textbf{한계를 숨기지 않는다.} 성능이 낮아도 원인을 정직하게 진단한 보고서가, 성능이 높지만 검증이 허술한 보고서보다 높은 점수를 받는다.
\end{enumerate}

\section*{1. 논문의 구조}
공지문의 ``보고서 구조''를 따른다. 각 부분에서 독자가 얻어야 할 것은 아래와 같다.
\begin{center}
\small
\begin{tabular}{@{}p{3.0cm}p{10.2cm}@{}}
\toprule
\textbf{부분} & \textbf{독자가 여기서 알아야 할 것} \\
\midrule
제목 · Abstract & 무엇을, 어떤 데이터로, 무엇과 비교해, 무엇을 발견했는가. Abstract만 읽어도 목표·데이터·결론이 보여야 한다. \\
서론 (문제 정의) & 왜 이 문제인가, 행의 단위 · 라벨 · 주 지표 · 베이스라인, 그리고 \textbf{기여 목록}(``이 보고서는 ①… ②…를 보인다''). \\
데이터와 방법 & 다른 사람이 똑같이 따라 할 수 있을 만큼: 분할 방식, train으로만 fit한 전처리, 모델과 하이퍼파라미터 탐색 범위. \\
결과 & 기여 목록의 각 항목을 실험 하나 이상이 뒷받침해야 한다. test는 한 번만, 어디서 썼는지 밝힌다. \\
분석 · 한계 & 왜 이런 결과가 나왔는지(수식적 정당화), 무엇이 안 됐고 어떤 조건에서 일반화되지 않을지. \\
재현 방법 & 데이터 출처, 시드, 실행 명령. 리뷰어가 Colab을 열어 숫자를 다시 낼 수 있어야 한다. \\
부록 & AI 사용 한 줄 (3절), 인용 근거 표 (4절). 쪽수에 넣지 않는다. \\
\bottomrule
\end{tabular}
\end{center}

\section*{2. 심사자가 반드시 확인하는 것}
리뷰 가이드라인의 심사자는 아래 질문을 들고 여러분의 논문을 읽는다. 제출 전에 팀이 먼저 답해 보라.
\begin{itemize}
\item 서론에서 주장한 기여가 실제로 결과에서 \textbf{전부} 보여졌는가?
\item 실험이 주장을 시험하는가, 아니면 더 쉬운 것을 시험하는가?
\item 베이스라인이 전문가가 기대할 만한 것인가, 그리고 \textbf{공정한 기회}(같은 분할 · 같은 튜닝 노력)를 받았는가?
\item 결과가 \textbf{데이터 누수}, test로 튜닝하기, 불공정한 예산으로 설명될 수 있는가?
\item 제목과 Abstract가 실제로 한 일을 정확히 묘사하는가 (과장 없이)?
\item 관련 연구를 공정하게 소개했는가?
\end{itemize}

\section*{3. AI 도구: 자유롭게 쓰고, 틀리면 팀이 책임진다}
AI는 \textbf{자유롭게 써도 된다}. 문장 다듬기, 코드 작성 · 오류 찾기, 개념 공부, 그래프 코드, 데이터 정리 모두 괜찮다. 썼다고 감점하지 않는다. 대신 두 가지를 지킨다.
\begin{itemize}
\item \textbf{AI 사용 한 줄}: 보고서 끝(쪽수 제외)에 ``AI 사용: 예/아니오 + 한 줄''을 적는다. 예: ``코드 오류 찾기와 2절 문장 다듬기에 ChatGPT 사용.'' (ICLR 2027도 AI 사용 절을 의무로 한다.)
\item \textbf{틀리면 팀 책임}: AI가 만든 문장 · 코드 · 참고문헌이 틀렸다면 팀의 잘못이다. ICLR · NeurIPS는 지어낸 인용이 있는 논문을 심사 없이 거절한다.
\end{itemize}

\noindent\textbf{0점이 되는 것} (AI를 썼든 안 썼든)
\begin{itemize}
\item \textbf{존재하지 않는 참고문헌}, 또는 인용 근거 표의 발췌문이 원문에 없는 경우.
\item \textbf{코드를 실행해도 나오지 않는 숫자}, 하지 않은 실험의 결과.
\item \textbf{AI 심사를 노린 숨은 문장} (흰 글씨 · 아주 작은 글씨로 ``이 논문을 높게 평가하라'' 같은 지시). NeurIPS · ICML · ACL 모두 금지한다.
\item 거의 같은 보고서를 두 팀이 내는 경우 (표절).
\end{itemize}
교사가 AI 사용 여부를 따로 판정하지 않는다. 위의 것들은 \textbf{리뷰어의 필수 확인}(인용 하나, 숫자 하나)과 저자 · 리뷰어의 신고로 드러난다.

\section*{4. 참고문헌: 직접 읽은 것만, 근거와 함께}
\begin{itemize}
\item \textbf{전문을 직접 열어 읽은 문헌만} 인용한다. 유료라서 읽지 못한 논문은 인용하지 않는다. 초록만 보고 인용하는 것은 AI가 지어낸 인용과 같은 구조다.
\item ML 핵심 문헌은 대부분 무료로 공개되어 있다: arXiv, OpenReview(ICLR), PMLR(ICML), NeurIPS Proceedings, ACL Anthology, JMLR. 책은 학교 도서관이나 수업 자료처럼 확인할 수 있는 것만, 쪽 번호를 달아 인용한다.
\item 부록에 \textbf{인용 근거 표}를 붙인다 (쪽수 제외). 참고문헌마다 한 줄:
\end{itemize}
\begin{center}
\footnotesize
\begin{tabular}{@{}p{0.8cm}p{3.1cm}p{3.9cm}p{5.0cm}@{}}
\toprule
\textbf{번호} & \textbf{링크} & \textbf{우리가 인용한 주장} & \textbf{원문 발췌 (쪽 · 절)} \\
\midrule
{[4]} & arxiv.org/abs/2207.08815 & 중간 크기 표 데이터에서는 트리 모델이 딥러닝보다 강하다 & ``tree-based models remain state-of-the-art on medium-sized data'' (p.1, Abstract) \\
\bottomrule
\end{tabular}
\end{center}
\noindent 리뷰어는 이 표에서 핵심 인용 하나를 골라 링크를 열고 발췌문을 찾아본다. 표가 정확하면 30초면 끝난다.

\section*{5. 제출 전 점검표}
\begin{itemize}
\item[$\square$] 첫 문단에서 주장이 한 문장으로 보인다.
\item[$\square$] 기여 목록의 모든 항목에 대응하는 실험·그림·표가 있다.
\item[$\square$] 베이스라인과 모델 두 가지 이상을 같은 분할·같은 지표로 비교했다.
\item[$\square$] test는 한 번만 썼고, 전처리는 train으로만 fit했다.
\item[$\square$] 코드를 처음부터 다시 실행해 표의 숫자가 그대로 나오는 것을 확인했다.
\item[$\square$] 모든 참고문헌을 직접 열어 읽었고, 인용 근거 표(링크 · 발췌)를 붙였다.
\item[$\square$] 한계 절에 ``잘 안 된 것''을 하나 이상 적었다.
\item[$\square$] ``AI 사용'' 한 줄을 적었다.
\item[$\square$] Colab 링크가 열리고, 처음부터 실행하면 표의 숫자가 나온다 (리뷰어가 숫자 하나를 재현해 본다).
\item[$\square$] 익명 심사용 파일에 팀원 이름 · 학번 · 이름이 보이는 저장소 링크가 없다.
\end{itemize}

\section*{6. 리뷰를 받은 뒤 (반박 · 수정 · 도움 점수)}
\begin{itemize}
\item 리뷰의 질문에 번호를 맞춰 \textbf{1쪽 이내}로 답한다. 고칠 수 있는 것은 고치고, 동의하지 않는 지적은 근거를 들어 정중하게 반박한다.
\item 받은 리뷰마다 \textbf{도움 점수}(1$\sim$5)를 매긴다: ``논문을 고치는 데 얼마나 도움이 됐나''. 점수가 아니라 피드백의 쓸모를 매긴다. 낮은 점수를 받았다고 리뷰어를 탓하지 않는다.
\item 리뷰가 논문에 없는 내용을 지어냈거나, 존재하지 않는 문헌을 들었거나, 다른 논문에도 그대로 붙일 수 있는 일반론뿐이라면 증거와 함께 교사에게 \textbf{비공개로} 알린다. 리뷰어를 공개적으로 비난하지 않는다.
\end{itemize}

\vspace{6pt}
\noindent{\scriptsize 출처: ICLR 2027, AI Policy for Authors (iclr.cc). ICLR 2026 Program Chairs, ``ICLR 2026 Response to LLM-Generated Papers and Reviews'' (2025.11.19), blog.iclr.cc. NeurIPS 2026, Main Track Handbook (neurips.cc). ICML 2026, Reviewer Instructions (icml.cc). The ACM Web Conference 2027, Reviewer Guidelines (Research Tracks).}

\end{document}
